\section{Conclusion}
\label{sec:conclusion}

We presented \method, a hierarchical language modeling framework that unifies dynamic semantic compression with balanced discrete reconstruction and generation through continuous--discrete co-adaptation. Dynamic compression and generation-oriented supervision encourage compact, predictive latent representations, while the decoder realizes latent semantics and handles residual uncertainty from context. Token-prediction gradients drive encoder--decoder co-adaptation, whereas latent-prior fitting remains decoupled from representation learning. Systematic component analyses and controlled evaluations on five downstream tasks, scaled to about 210B training tokens, support these design principles and show competitive, unsaturated scaling. In a separate resource-constrained run, the final checkpoint trained on 346.03B source tokens achieves a 60.10 average over eight tasks and surpasses Llama-3-8B on HellaSwag and StoryCloze, although that baseline reports over $43{\times}$ more pretraining tokens. These results indicate successful transfer of the Qwen-initialized causal decoder's generation capability and substantial scaling headroom, especially because the latent DiT remains far from convergence. More broadly, \method shows that continuous-latent hierarchical modeling can integrate effectively with established token-space generators, including autoregressive and masked-diffusion decoders, rather than depending on a single decoder paradigm.
